\documentclass{ctexbook}
\begin{document}

\chapter{第二章 过程分析与建模基础}
第二章引言属于概述式引言。近年来智能制造和工业互联网在国民经济中蓬勃发展并起到关键作用。
针对上述现实需求，本章首先给出总体框架，其次设计分析流程，最后完成工况分析。

\section{2.1 工艺流程与变量分析}
本节分析实际工业工艺流程中的主要测量变量与输入输出关系。

\chapter{第三章 自适应动态状态估计方法}
第2章给出了基线模型，但在复杂工况下存在泛化漂移瓶颈。
近年来智能制造和工业互联网飞速发展，在国民经济和战略需求中发挥着核心支撑作用，具有广泛应用前景。
针对上述问题，本章提出自适应重加权机制。
本章组织如下：3.1 节介绍模型结构，3.2 节介绍更新准则，3.3 节给出实验验证。同时，本章首先建立数学模型，其次推导更新算法，最后完成实验对比。

\section{3.1 状态估计系统框架}
第2章给出了基线模型，但在复杂工况下存在泛化漂移瓶颈。针对上述问题，本章提出自适应重加权机制，本节首先建立数学模型，其次推导更新算法，最后完成实验验证，并针对泛化漂移瓶颈与自适应重加权展开深入研究。

\subsection{3.1.1 特征投影算子构建}
在构建特征投影算子时，现有研究通常面临多重挑战：首先是高维观测带来的特征冗余瓶颈，其次是多模态采样不同步导致的分布漂移难题。为了克服上述全部挑战，本小节展开深入推导。
针对上节输出的异步特征表征，本小节构建时间连续投影算子。

\subsection{3.1.2 投影矩阵在线更新}
本小节建立投影矩阵在线更新准则。给定输入序列向量，通过梯度投影更新参数。
\begin{equation}
\mathbf{y} = \mathbf{W} \mathbf{x} + \mathbf{b}
\end{equation}
上式中，首先计算矩阵相乘，再将乘积进行矩阵转置并加上偏置，最后应用激活函数计算出各项最终输出概率。

\section{3.2 实验验证与对比分析}
本节重点给出实验对比结果。下文 3.2.1 节介绍测试平台设置，3.2.2 节给出各算法对比指标与统计显著性分析。

\subsection{3.2.1 实验平台与基准模型}
本小节交代实验验证平台的数据集划分与评价指标。给定输入序列，模型依次完成测试。

\subsection{3.2.2 评价指标与对比结果}
本小节给出对比实验中的投影矩阵计算方法与指标分析。
\begin{equation}
\mathbf{z} = \mathbf{A} \mathbf{y}
\end{equation}
式中：$\mathbf{A}$ 为正交投影矩阵，$\mathbf{y}$ 为输入时序特征向量。该变换将观测向量映射到主成分补空间，用于消除稳态噪声干扰。

\section{3.3 本章小结}
本章针对非平稳工况下的状态估计漂移瓶颈，提出了自适应动态重加权方法。
如图 \ref{fig:framework} 与表 \ref{tab:comparison} 所示，本文方法显著降低了估计误差。
在此基础上，根据文献 \cite{smith2024} 的收敛准则完成训练：
\begin{equation}
\mathcal{L} = \alpha \mathcal{L}_1 + \beta \mathcal{L}_2
\end{equation}

\end{document}
